% ================== ЛАБОРАТОРНАЯ РАБОТА №2 ==================
\documentclass[14pt,a4paper]{extarticle}

\usepackage[utf8]{inputenc}
\usepackage[russian]{babel}
\usepackage{geometry}
\geometry{left=30mm,right=15mm,top=20mm,bottom=20mm}

\usepackage{indentfirst}
\setlength{\parindent}{1.25cm}

\usepackage{amsmath,amssymb}
\usepackage{graphicx}
\usepackage{float}
\usepackage{xcolor}
\usepackage{hyperref}
\usepackage{caption}
\usepackage{setspace}
\setstretch{1.0}

\usepackage{titlesec}
\titleformat{\section}{\normalfont\normalsize\bfseries}{\thesection}{1em}{\MakeUppercase}
\titleformat{\subsection}{\normalfont\normalsize\bfseries}{\thesubsection}{1em}{}
\titlespacing{\section}{\parindent}{3.5ex plus 1ex minus .2ex}{2.3ex plus .2ex}
\titlespacing{\subsection}{\parindent}{3.25ex plus 1ex minus .2ex}{1.5ex plus .2ex}

\usepackage{fancyhdr}
\pagestyle{fancy}
\fancyhf{}
\fancyfoot[R]{\thepage}
\renewcommand{\headrulewidth}{0pt}
\fancypagestyle{plain}{
  \fancyhf{}
  \fancyfoot[R]{\thepage}
  \renewcommand{\headrulewidth}{0pt}
}

\renewcommand{\contentsname}{\centering СОДЕРЖАНИЕ}

\begin{document}

\thispagestyle{empty}

\begin{center}
МИНИСТЕРСТВО ОБРАЗОВАНИЯ РЕСПУБЛИКИ БЕЛАРУСЬ\\[0.5em]
Учреждение образования\\[0.5em]
\textbf{<<Белорусский государственный университет информатики и радиоэлектроники>>}\\[2em]

Факультет компьютерных систем и сетей\\[0.5em]
Кафедра программного обеспечения информационных технологий
\end{center}

\vspace{3em}

\begin{center}
\textbf{ОТЧЁТ}\\[0.5em]
по лабораторной работе №2\\[0.5em]
по дисциплине\\[0.5em]
\textbf{<<Генетические и эволюционные вычисления>>}

\textbf{Тема: <<Прогнозирование значений процесса с помощью линейной нейронной сети>>}
\end{center}

\vspace{4em}

\hfill\begin{minipage}{0.55\textwidth}
Выполнил:\\
магистрант 1 курса магистратуры\\
факультета компьютерных систем и сетей\\
группы 556301\\[0.3em]
\textbf{Лебедевич Артём Владимирович}\\[1.2em]
Преподаватель:\\
доцент кафедры ПОИТ, канд. техн. наук\\
\textbf{Нестеренков Сергей Николаевич}
\end{minipage}

\vfill

\begin{center}
Минск, 2026~г.
\end{center}

\newpage

\tableofcontents

\newpage

\section{Цель работы}

Изучить принципы построения и обучения линейных нейронных сетей для задачи прогнозирования временных рядов. Реализовать линейную нейронную сеть, прогнозирующую следующее значение дискретного сигнала по $k$ предыдущим значениям, и исследовать влияние параметров на качество прогноза.

\section{Теоретические сведения}

\subsection{Линейная нейронная сеть}

Линейная нейронная сеть (ADALINE) реализует линейную комбинацию входных сигналов:

\begin{equation}
y = \sum_{i=1}^{k} w_i x_i + b = \mathbf{w}^T \mathbf{x} + b,
\end{equation}

где $\mathbf{w} = (w_1, \ldots, w_k)^T$~--- вектор весов, $b$~--- смещение, $\mathbf{x} = (x_1, \ldots, x_k)^T$~--- входной вектор.

\subsection{Метод наименьших квадратов}

Для оптимального подбора весов решается задача минимизации суммы квадратов ошибок. Аналитическое решение через псевдообратную матрицу:

\begin{equation}
\mathbf{W} = \mathbf{T} \cdot \mathbf{P}^+,
\end{equation}

где $\mathbf{P}^+ = \mathbf{P}^T (\mathbf{P} \mathbf{P}^T)^{-1}$~--- псевдообратная матрица Мура---Пенроуза, $\mathbf{T}$~--- вектор целевых значений.

\subsection{Формирование матрицы входов}

Для прогнозирования временного ряда матрица входов $P$ формируется по принципу скользящего окна: каждый столбец содержит $k$ последовательных значений сигнала, а соответствующий целевой выход~--- следующее за ними значение:

\begin{equation}
P = \begin{pmatrix}
x_1 & x_2 & \cdots & x_{N-k} \\
x_2 & x_3 & \cdots & x_{N-k+1} \\
\vdots & \vdots & \ddots & \vdots \\
x_k & x_{k+1} & \cdots & x_{N-1}
\end{pmatrix}, \quad
T = (x_{k+1}, x_{k+2}, \ldots, x_N).
\end{equation}

\section{Ход работы}

\subsection{Генерация сигнала}

Исследуемый сигнал задан формулой:
\begin{equation}
x(t) = \sin(2t), \quad t \in [0, 4\pi], \quad \Delta t = 0.1~\text{с}.
\end{equation}

Сгенерирован дискретный сигнал, содержащий 126 отсчётов.

\subsection{Построение линейной нейронной сети}

Реализована функция формирования матрицы входов $P$ по методу скользящего окна с параметром $k$ (число предыдущих значений). Веса нейронной сети определены аналитически через псевдообратную матрицу с использованием функции \texttt{numpy.linalg.pinv}.

\subsection{Моделирование}

При $k = 5$ линейная сеть демонстрирует высокую точность прогноза: предсказанный сигнал практически совпадает с эталонным. График ошибки показывает малые отклонения, сосредоточенные в начале сигнала (переходный процесс).

\subsection{Анализ влияния параметра $k$}

Исследовано влияние размера окна $k$ на качество прогноза для $k \in \{2, 3, 4, 5, 6, 7, 8, 9, 10\}$.

\begin{itemize}
\item При $k = 2$: MSE относительно высока~--- двух предыдущих значений недостаточно для точного прогноза.
\item При $k = 3 \ldots 5$: существенное снижение ошибки, сеть получает достаточно информации о структуре сигнала.
\item При $k \geq 5$: дальнейшее увеличение $k$ практически не улучшает результат, MSE выходит на плато.
\end{itemize}

\subsection{Анализ влияния частоты сигнала}

Исследовано влияние частотного параметра $\omega$ при фиксированном $k = 5$ для $\omega \in [0.5, 10.0]$.

С увеличением $\omega$ частота сигнала растёт, и при фиксированном шаге дискретизации $\Delta t = 0.1$ на один период приходится меньше отсчётов. Это приводит к росту ошибки прогнозирования, особенно при $\omega > 5$, когда теорема Котельникова (Найквиста---Шеннона) нарушается.

\section{Результаты}

\begin{enumerate}
\item Линейная нейронная сеть на основе псевдообратной матрицы обеспечивает высокую точность прогноза для гладких периодических сигналов.
\item Оптимальный размер окна для $x(t) = \sin(2t)$ при $\Delta t = 0.1$ составляет $k \approx 5$: дальнейшее увеличение не даёт значимого улучшения.
\item Качество прогноза деградирует при увеличении частоты сигнала из-за недостаточной дискретизации.
\end{enumerate}

\section{Выводы}

В ходе лабораторной работы реализована линейная нейронная сеть для прогнозирования временного ряда. Аналитическое решение через псевдообратную матрицу обеспечивает оптимальные веса в смысле наименьших квадратов. Экспериментально подтверждено, что качество прогноза зависит от размера окна $k$ и частоты дискретизации: необходим баланс между информативностью входного вектора и сложностью задачи.

\end{document}
